In this section, we explore a MIP formulation of the pricing problem \eqref{eq:RP_set_formulation} and discuss challenges in solving it.
Recall the reduced cost \eqref{eq:RM-redcost}.
Let $\bm{z}\in \{0,1\}^{\L}$ again be integer variables representing which clauses are used in the rule.
Consider the MIP
\begin{subequations}
    \label{eq:RP}
    \begin{align}
        \min \quad & \sum_{i \in \cN} \gamma_i - \sum_{i \in \cP} \mu_i \gamma_i +
        \lambda(\bm{z}) \label{eq:RP-obj} \\
        \st \quad & \gamma_i \leq 2 - (2 - \gamma_{il})z_l, & l \in \L, &\; i \in \cP, \label{eq:RP-pos} \\
        & \gamma_i \geq \gamma_{il} - \sum_{l' \in \L \fw \gamma_{il'} < \gamma_{il}} (\gamma_{il} - \gamma_{il'}) z_{l'} , & l \in \cL, &\; i \in \cN, \label{eq:RP-neg} \\
        & \gamma_i \geq 2 - \sum_{l \in \L} 2z_l, & l \in \L, \label{eq:RP-fix} \\
        & \sum_{l \in \L} z_l \leq D, \label{eq:RP-complex}\\
        & \bm{z}\in \{0,1\}^{\L},~ \bm{\gamma}\in [0,2]^\I.\label{eq:RP-vars}\qquad
    \end{align}
\end{subequations}
We also impose the one-clause-per-feature constraint as given in \eqref{eq:TP_one-per-feature}.

In this formulation, constraints \eqref{eq:RP-pos} ensure that $\gamma_i$ for positive $i$ reflects $\gamma_{ik},$
where $k$ is the chosen rule.
Strictly speaking, this is only true for $i \in \P$ for which $\mu_i > 0$.
For $i \in \P$ such that $\mu_i = 0$, the variable $\gamma_i$ is effectively redundant,
so these variables need not be considered.
For negative $i$, constraints \eqref{eq:RP-neg} and \eqref{eq:RP-fix} ensure this.
We will now prove the claim for both positive and negative points.

Denote by $\gopt, \zopt$ an optimal solution for this IP\@.
Recall that $\zopt$ is binary.
Define
\[\sopt = \{l \in \L \fw \zopt[l] = 1\},\]
and recall the definition of $f$ in \eqref{eq:f_min}.

\begin{claim}\label{claim:TP_ramp-consP}
    For all $i \in \P$ such that $\mu_i > 0,$ it holds that
    \[\gopt[i] = f_i(\sopt).\]
\end{claim}

\begin{proof}
    Let $i \in \P$ such that $\mu_i > 0$ and let $l \in \L$.
    If $\zopt[l] = 0$, the corresponding constraint \eqref{eq:RP-pos} becomes redundant.
    If $\zopt[l] = 1,$ the constraint becomes $\gamma_i \leq \gamma_{il}.$
    Therefore, if $\sopt \neq \emptyset$ it holds that
    \[\gamma_i \leq \umin{l \in \sopt} \gamma_{il}.\]
    If $\sopt = \emptyset$, we have that $\gamma_i \leq 2$ by the variable bound.
    In any case,
    \[
        \gamma_i \leq f_i(\sopt).
    \]
    There are no other constraints involving $\gamma_i$.
    Since $\gamma_i$ has a negative coefficient in the objective, this inequality will be an equality in an optimal solution.
    \qedhere
\end{proof}

\begin{claim}\label{claim:TP_ramp-consN}
    For all $i \in \N$ it holds that
    \[\gopt[i] = f_i(\sopt).\]
\end{claim}


\begin{proof}
    Let $i \in \N$.
    If $\sopt = \emptyset,$ then $\gopt[i] \geq 2 = f_i(\emptyset)$ by constraint \eqref{eq:RP-fix}.
    So if $\sopt = \emptyset$, the statement holds.

    Assume $\sopt \neq \emptyset.$
    Observe that constraint $\eqref{eq:RP-fix}$ becomes redundant, as there is at least one $l$ such that $z_l = 1$.
    Let $l^* \in \sopt$ such that \[\gamma_{il^*} = \umin{l \in \sopt} \gamma_{il}.\]
    For this $i$ and $l$, the constraint \eqref{eq:RP-neg} is
    \[\gamma_i \geq \gamma_{il}.\]
    Now let $l \in \L$.
    We consider the right hand side of \eqref{eq:RP-neg}.
    We have that
    \[\begin{split}
          \gamma_{il} - \sum_{l' \in \L\,:\,\gamma_{il'} < \gamma_{il}} (\gamma_{il} - \gamma_{il'}) z_l
          &= \gamma_{il} - \sum_{l' \in \sopt \,:\,\gamma_{il'} < \gamma_{il}} (\gamma_{il} - \gamma_{il'}).
    \end{split}\]
    If the sum is empty, then it must hold that $\gamma_{il} = \gamma_{il^*}$ and the right hand side equals $\gamma_{il^*}$.
    If the sum is non-empty, we have
    \[
        \begin{split}
            \gamma_{il} - \sum_{l' \in \sopt \,:\,\gamma_{il'} < \gamma_{il}} (\gamma_{il} - \gamma_{il'}) &\leq \gamma_{il} - \umax{l' \in \sopt \,:\,\gamma_{il'} < \gamma_{il}} (\gamma_{il} - \gamma_{il'}) \\
            &= \umin{l' \in \sopt \,:\,\gamma_{il'} < \gamma_{il}} \gamma_{il'} \\
            &= \gamma_{il^*}.
        \end{split}
    \]
    Hence, $\gamma_i \geq \gamma_{il^*} = \umin{l \in \sopt} \gamma_{il}$ is the tightest constraint on $\gamma_i$, induced by $l^*$.

    Combining the case $\sopt = \emptyset$ and $\sopt \neq \emptyset,$ we get that
    \[
        y_i \geq f_i(\sopt),
    \]
    is the tightest lower-bound on $\gamma_i$.
    Since $\gamma_i$ has a positive coefficient in the objective, this inequality will be an equality for an optimal solution.
    From this, the result follows.
\end{proof}

\begin{note}
    In the process of constructing and implementing this MIP, constraint \eqref{eq:RP-fix} was overlooked.
    Hence, it was not implemented when the experiments in this report were conducted.
    Luckily, this mistake does not affect the experiments for various reasons.
    First, the solving the MIP could not be done in reasonable time for the PD-L1 experiments,
    hence only the unaffected heuristic pricer was used.
    If the exact pricer implementation without constraint \eqref{eq:RP-fix} is used anyway, the effect is still limited.
    We note that for any non-zero solution $\bm{z} \in \{0,1\}^\L,$ the variables $\bm{\gamma}$ are correctly modeled.
    Hence, the reduced cost can only be incorrect when $\bm{z} = 0$.
    For $\bm{z} = 0$, it can happen that some $\gamma_i$ are too low for negative $i$.
    If so, the reduced cost is underestimated.
    This is only problematic if, due to this error, the empty rule becomes the rule with the lowest reduced cost.
    In this case, the rule would not be added to the master problem in our implementation,
    and pricing would continue (with new shadow prices),
    which may cause a rule with true negative reduced cost to be found.
    However, the problem of underestimating $\gamma_i$ for negative $i$ only occurs if there is no
    clause $l$ such that $\gamma_{il} = 2$.
    This is the case when there is no clause that covers (further than the margin) the negative point.
    In practice, we do not expect this to happen often,
    as we include upper-bounds with low thresholds and lower-bounds with high thresholds.
\end{note}

The LP relaxation of \eqref{eq:RP} is not tight.
The LP relaxation allows for \emph{spreading}, by which several entries of $\bm{z}$ admit a fractional value.
\begin{example}[Spreading]\label{ex:branch1}
    \begin{figure}
        \includesvg[width=0.5\textwidth]{gen/branch_ex1.svg}
        \caption{
            An example problem with a two-dimensional feature space, composed of one positive datapoint $\bm{X}_2 = (2\tfrac{1}{2}, 2\tfrac{1}{2})$, two negative datapoints $\bm{X}_1 = (1\tfrac{1}{2}, 1\tfrac{1}{2}), \bm{X}_3 = (3\tfrac{1}{4}, 3\tfrac{1}{4})$ and clauses $a$, $b$, $f$ representing $X_{\cdot 1} \geq 3$,$X_{\cdot 2} \geq 3$ $X_{\cdot 1} \geq \tfrac{1}{4}$, respectively.
            The margin width $\beta=1.$
            The boundary of the margin is indicated with dotted lines.}
        \label{fig:branch1}
    \end{figure}
    Let
    %
    \[
        \mathcal{N} = \{1, 3\}, \quad \mathcal{P} = \{2\}, \quad
        \mathcal{L} = \{a, b, f\}.
    \]
    Consider the two-dimensional feature space and lower-bounds in \autoref{fig:branch1}.
    The matrix $(\gamma_{il})$ is given by
    %
    \[
        \begin{array}{c|ccc}
            & a    & b    & f \\
            \hline
            1 & 0    & 0    & 1 \\
            2 & 1/2  & 1/2  & 1 \\
            3 & 1/4  & 1/4  & 1 \\
        \end{array}
    \]
    Choose $D = 3.$ In the initial pricing round, $\mu_i = 1$ and $\lambda = 0$.
    The LP relaxation of~\eqref{eq:RP} attains an optimal objective value of $-1/2,$
    by
    %
    \[
        \bm{z} = \bigl(\tfrac{1}{2},\, \tfrac{1}{2},\, 0\bigr), \qquad
        \gamma_1 = 0, \quad \gamma_2 = \tfrac{3}{4}, \quad \gamma_3 = \tfrac{1}{4}.
    \]
    Given $\bm{z}$, the values of $\bm{\gamma}$ can be computed directly.
    The optimal objective value of the IP is $-1/4, $ attained by several solutions, examples being
    \[\bm{z}=(1,0,0), \quad \bm{z}=(0,1,0), \quad \bm{z}=(1,1,0).\]
    For any of these solutions, it holds that $\gamma_2 \leq 1/2$ because at least one of $a$ or $b$ is used.
    Meanwhile, we have $\gamma_2 = \frac{3}{4}$ in an optimal LP solution, like the one shown above.
    We refer to this as spreading, meaning the use of fractional values in $\bm{z}$ to admit larger $\gamma_i$ for $i \in \P,$ leading to a lower objective value.

    In this example, the issue is resolved when branching on $z_a$ or $z_b$.
    Under any of the branching decisions
    \[z_a = 0,\; z_a = 1,\; z_b = 0, \text{ or } z_b = 1,\]
    the resulting LP solution is integer and $\gamma_2 = 1/2.$
\end{example}

%\begin{example}[Spreading and branching]\label{ex:branch2}
%    \begin{figure}
%        \includesvg[width=0.5\textwidth]{gen/branch_ex2.svg}
%        \caption{
%            An example problem with a two-dimensional feature space, composed of one positive datapoint $\bm{X}_2 = (2\tfrac{1}{2}, 2\tfrac{1}{2})$, two negative datapoints $\bm{X}_1 = (1\tfrac{1}{2}, 1\tfrac{1}{2}), \bm{X}_3 = (3\tfrac{1}{4}, 3\tfrac{1}{4})$ and clauses $a, b, c, d, f$, representing $X_{\cdot 1} \geq 3$, $X_{\cdot 2} \geq 3$, $X_{\cdot 2} \geq 2\tfrac{3}{4}$, $X_{\cdot 2} \geq 2\tfrac{3}{4}$ and $X_{\cdot 1} \geq \tfrac{1}{4}$, respectively.
%            The margin width $\beta=1.$
%            The boundary of the margin is not shown.}
%        \label{fig:branch2}
%    \end{figure}
%
%    We consider the problem discussed in \autoref{ex:branch1} with two additional clauses $c$ and $d$, see \autoref{fig:branch2}.
%    The matrix $(\gamma_{il})$ is given by
%    %
%    \[
%        \begin{array}{c|ccccc}
%              & a    & b    & c    & d    & f \\
%            \hline
%            1 & 0    & 0    & 0    & 0    & 1 \\
%            2 & 1/2  & 1/2  & 3/4  & 3/4  & 1 \\
%            3 & 1/4  & 1/4  & 1/2  & 1/2  & 1 \\
%        \end{array}
%    \]
%    The LP relaxation of~\eqref{eq:RP} attains an optimal objective value of $-1/2,$
%    achieved for example by
%    %
%    \[
%        \bm{z} = \bigl(\tfrac{1}{4},\, \tfrac{1}{4},\, \tfrac{1}{4},\,
%        \tfrac{1}{4},\, 0\bigr), \qquad
%        \gamma_1 = 0, \quad \gamma_2 = \tfrac{7}{8}, \quad \gamma_3 = \tfrac{3}{8}.
%    \]
%    %
%    In any solution to the IP that uses clause $a$ or $b$, it holds that $\gamma_2 \leq 1/2.$
%    In any solution that uses clause $c$ or $d$, it holds that $\gamma_2 \leq 3/4.$
%    Spreading is not necessarily prevented by branching on a variable with fractional value.
%    Suppose we branch on $z_d.$ If $z_d = 0$ is fixed, there exists optimal LP solution
%    \[
%        \bm{z} = \bigl(\tfrac{1}{4},\, \tfrac{1}{4},\, \tfrac{1}{4},\,
%        \tfrac{1}{2},\, 0\bigr), \qquad
%        \gamma_1 = 0, \quad \gamma_2 = \tfrac{7}{8}, \quad \gamma_3 = \tfrac{3}{8},
%    \]
%    with objective value $-1/2.$
%    If $z_d = 1$ is fixed, an optimal solution is given by
%    \[
%        \bm{z} = \bigl(\tfrac{1}{2},\, \tfrac{1}{2},\, \tfrac{1}{4},\,0,\, 1\bigr), \qquad
%        \gamma_1 = 0, \quad \gamma_2 = \tfrac{3}{4}, \quad \gamma_3 = \tfrac{1}{4},
%    \]
%    again with objective value $-1/2.$
%    If constraint \eqref{eq:TP_one-per-feature} is also added, an optimal solution when $z_d = 1$ is fixed is given by
%    \[
%        \bm{z} = \bigl(\tfrac{1}{2},\, 0,\, \tfrac{1}{2},\,1,\, 1\bigr), \qquad
%        \gamma_1 = 0, \quad \gamma_2 = \tfrac{3}{4}, \quad \gamma_3 = \tfrac{3}{8},
%    \]
%    and has objective value $-3/8.$
%    In all cases, an optimal LP solution after branching is still fractional.
%    An analogous result is observed when branching on $z_c$.
%    Branching on $z_a$ or $z_b$ is better.
%    If $z_a = 1$ is fixed, for example, an optimal LP solution has objective value $-1/4$, the optimal IP objective value.
%    The solution found by SCIP is integer, namely the solution that uses clauses $a$ and $b$.
%    If $z_a = 0$ is fixed, the optimal LP objective value is $-9/20$, attained for example by
%    \[
%        \bm{z} = \bigl(0,\, \tfrac{1}{5},\, \tfrac{2}{5},\,
%        \tfrac{2}{5},\, 0\bigr), \qquad
%        \gamma_1 = 0, \quad \gamma_2 = \tfrac{9}{10}, \quad \gamma_3 = \tfrac{9}{20}.
%    \]
%    \todo{How would SCIP choose a variable to branch on? Would it systematically avoid $z_d$ and $z_c$?
%    How would it work in more complicated examples? Is it then possible to choose these 'useless' branching variables?}
%
%    An alternative approach is branching on a set of variables.
%    For instance, choose $S=\{a, b\}$ and consider the two branches
%    \[\sum_{l \in S} z_l \geq 1, \text{ and } \sum_{l \in S} z_l = 0.\]
%    This branching strategy by itself is not more effective than variable branching.
%    However, each branch induces problem-specific variable bounds not captured by the LP\@.
%    In the former branch, at least one of $a$ or $b$ is chosen, and therefore \[\gamma_2 \leq 1/2.\]
%    Adding this constraint, which is equivalent to an updated upper bound on $\gamma_2$, this branch yields an integer solution.
%    In the latter branch, it follows that \[\gamma_3 \geq 1/2.\]
%    Even with this induced constraint, the branch does not yield an integer solution.
%    Instead, spreading is observed over $c$ and $d$.
%    Note that the set branching is a generalization of variable branching, as branching on some $S_1$ with one element is equivalent to branching on a variable.
%\end{example}